\documentclass[tikz,border=6pt]{standalone}
\usepackage{ctex}
\usepackage{amsmath}
\usetikzlibrary{calc, arrows.meta}

\begin{document}
\begin{tikzpicture}[scale=1.15, thick]

% 椭圆 x^2/4 + y^2 = 1，a=2, b=1
% 直线 y = k(x-1)，取 k=0.8 为示意
% 韦达定理示意：A(x1,y1), B(x2,y2) 两交点，标注 x1+x2, x1*x2

\def\a{2.0}
\def\b{1.0}

% 坐标轴
\draw[->, gray, thin] (-2.8, 0) -- (2.8, 0) node[right, font=\footnotesize] {$x$};
\draw[->, gray, thin] (0, -1.5) -- (0, 1.8) node[above, font=\footnotesize] {$y$};
\node[font=\footnotesize, below left] at (0,0) {$O$};

% 椭圆
\draw[black, thick] (0,0) ellipse [x radius=\a, y radius=\b];

% 顶点标注
\node[font=\footnotesize, below] at (2,0) {$a$};
\node[font=\footnotesize, below] at (-2,0) {$-a$};
\node[font=\footnotesize, right] at (0,1) {$b$};

% 直线 y = 0.8(x-1)，过点(1,0)，斜率 k=0.8
% 代入椭圆: x^2/4 + 0.64(x-1)^2 = 1
% x^2/4 + 0.64x^2 - 1.28x + 0.64 = 1
% 0.25x^2 + 0.64x^2 - 1.28x + 0.64 - 1 = 0
% 0.89x^2 - 1.28x - 0.36 = 0
% x = (1.28 ± sqrt(1.6384 + 4*0.89*0.36)) / (2*0.89)
% x = (1.28 ± sqrt(1.6384 + 1.2816)) / 1.78
% x = (1.28 ± sqrt(2.92)) / 1.78
% sqrt(2.92) ≈ 1.7088
% x1 = (1.28 + 1.7088)/1.78 ≈ 1.679
% x2 = (1.28 - 1.7088)/1.78 ≈ -0.241
% y1 = 0.8*(1.679-1) = 0.543
% y2 = 0.8*(-0.241-1) = -0.993

\coordinate (A) at (1.679, 0.543);
\coordinate (B) at (-0.241, -0.993);
\coordinate (M) at ({(1.679-0.241)/2}, {(0.543-0.993)/2}); % 中点

% 过点 P(1,0)
\coordinate (P) at (1, 0);
\fill[red!60] (P) circle [radius=1.8pt];
\node[font=\footnotesize, below] at (P) {$P$};

% 直线（延长一点）
\draw[blue, thick] (-0.7, {0.8*(-0.7-1)}) -- (2.1, {0.8*(2.1-1)});

% 交点 A, B
\fill[black] (A) circle [radius=2pt];
\fill[black] (B) circle [radius=2pt];
\node[font=\footnotesize, above right] at (A) {$A(x_1,y_1)$};
\node[font=\footnotesize, below left] at (B) {$B(x_2,y_2)$};

% 中点 M
\fill[orange] (M) circle [radius=2pt];
\node[font=\footnotesize, above right, orange] at (M) {$M$（弦中点）};

% 韦达定理示意：x轴上标注 x2, x1
\draw[dashed, gray] (A) -- (1.679, 0);
\draw[dashed, gray] (B) -- (-0.241, 0);

% x1, x2 标注
\node[font=\footnotesize, below, red!70] at (1.679, 0) {$x_1$};
\node[font=\footnotesize, below, red!70] at (-0.241, 0) {$x_2$};

% 双箭头标注 x1+x2（在 x 轴下方）
\draw[{Stealth[length=5pt]}-{Stealth[length=5pt]}, red!60, thin]
    (-0.241, -0.22) -- (1.679, -0.22)
    node[midway, below, font=\footnotesize, red!60] {$x_1{-}x_2$（差）};

% 标注韦达定理公式（右侧）
\node[align=left, font=\small] at (0.3, 1.5)
  {韦达定理：};
\node[align=left, font=\footnotesize] at (1.5, 1.25)
  {$x_1{+}x_2 = -\dfrac{q}{p}$};
\node[align=left, font=\footnotesize] at (1.5, 0.85)
  {$x_1 x_2 = \dfrac{r}{p}$};

% 弦长公式
\node[font=\small, align=center] at (0, -1.45)
  {$|AB| = \sqrt{1{+}k^2}\cdot\sqrt{(x_1{+}x_2)^2 - 4x_1 x_2}$};

\end{tikzpicture}
\end{document}
